\section{Positive-diagonal Matignon stability and orbit reduction}
\label{sec:framework}

For the commensurate Caputo system
\[
{}^C D_t^\alpha x=Ax,\qquad 0<\alpha<1,
\]
asymptotic stability is equivalent to
\[
|\arg\lambda|>\alpha\pi/2
\quad\forall\lambda\in\sigma(A),
\]
with zero excluded \cite{Matignon1996,BrandiburGarrappaKaslik2021}. Define
\[
\Sigma_\alpha=\{z\neq0:|\arg z|>\alpha\pi/2\}.
\]

\begin{definition}[Positive-diagonal Matignon stability]
\[
\mathcal F_\alpha^{(n)}
=
\{A:\sigma(DA)\subset\Sigma_\alpha\ \forall D\succ0\text{ diagonal}\},
\]
\[
\mathcal D_H^{(n)}
=
\{A:DA\text{ Hurwitz }\forall D\succ0\text{ diagonal}\},
\qquad
\mathcal P_\alpha^{(n)}
=
\mathcal F_\alpha^{(n)}\setminus\mathcal D_H^{(n)}.
\]
\end{definition}

All interiors are taken in the Euclidean topology of \(\mathbb R^{n\times n}\). Since
\[
DA=D(AD)D^{-1},
\]
left and right positive-diagonal multiplication are spectrally equivalent. A common positive scaling of \(D\) multiplies all eigenvalues by a positive scalar and therefore does not change their arguments. Thus the diagonal orbit is naturally projective. This is a special case of generalized region D-stability \cite{Kushel2019,KushelPavani2022}; our task is to eliminate the diagonal multiplier explicitly in low dimension.

Assume \(-A\) is a strict P-matrix and define signed principal minors
\[
m_I=(-1)^{|I|}\det A[I]>0,
\qquad
p_i=m_{\{i\}}=-a_{ii}>0,
\]
and normalized minors
\[
\beta_I=\frac{m_I}{\prod_{i\in I}p_i},
\qquad |I|\ge2.
\]

\begin{lemma}[Positive-diagonal orbit reduction]
\label{lem:simplex-reduction}
Let \(D=\operatorname{diag}(d_1,\dots,d_n)\succ0\) and
\[
\det(\lambda I-DA)
=
\lambda^n+c_1(D)\lambda^{n-1}+\cdots+c_n(D).
\]
Then
\[
c_k(D)=\sum_{|I|=k}m_I\prod_{i\in I}d_i.
\]
With
\[
s=c_1(D)=\sum_i p_i d_i,
\qquad
x_i=\frac{p_i d_i}{s},
\]
one has \(x\in\Delta_{n-1}^\circ\) and
\[
\frac{c_k(D)}{s^k}
=
B_k(x)
=
\sum_{|I|=k}\beta_I\prod_{i\in I}x_i.
\]
Conversely every \(x\in\Delta_{n-1}^\circ\) is obtained from a positive diagonal \(D\), uniquely up to common positive scale.
\end{lemma}

\begin{proof}
The coefficient of \(\lambda^{n-k}\) is
\[
(-1)^k\sum_{|I|=k}\det(DA[I])
=
\sum_{|I|=k}m_I\prod_{i\in I}d_i,
\]
because \(D\) is diagonal. Substituting \(d_i=sx_i/p_i\) gives the normalized formula. Conversely, for any \(x\in\Delta_{n-1}^\circ\) and \(s>0\), the choice \(d_i=sx_i/p_i\) reconstructs \(D\).
\end{proof}

After \(\lambda=sz\), the normalized characteristic polynomial is
\[
z^n+z^{n-1}+B_2(x)z^{n-2}+\cdots+B_n(x)=0.
\]
Thus the complete positive-diagonal orbit modulo scale is exactly the open simplex.

For \(n=3\), write
\[
m_{ij}=\det A[\{i,j\}],
\qquad
q=-\det A,
\]
and define the four orbit invariants
\[
\beta_{12}=\frac{m_{12}}{p_1p_2},
\quad
\beta_{13}=\frac{m_{13}}{p_1p_3},
\quad
\beta_{23}=\frac{m_{23}}{p_2p_3},
\quad
\kappa=\frac{q}{p_1p_2p_3}.
\]
Positive left-diagonal scaling multiplies numerator and denominator of each ratio by the same row factors, so all four are invariant.

\begin{figure}[t]
\centering
\begin{minipage}[t]{0.48\linewidth}\centering
\includegraphics[width=\linewidth]{panel01_spectral_geometry.pdf}
\end{minipage}\hfill
\begin{minipage}[t]{0.48\linewidth}\centering
\includegraphics[width=\linewidth]{panel02_dimension_threshold.pdf}
\end{minipage}

\begin{minipage}[t]{0.48\linewidth}\centering
\includegraphics[width=\linewidth]{panel03_simplex_objective.pdf}
\end{minipage}\hfill
\begin{minipage}[t]{0.48\linewidth}\centering
\includegraphics[width=\linewidth]{panel04_cyclic_band.pdf}
\end{minipage}
\caption{Structural geometry of the problem. (a) The Matignon sector contains a right-half-plane sliver absent from Hurwitz stability. (b) The robust fractional-only class first acquires nonempty interior in dimension three. (c) A representative normalized simplex objective; Theorem~\ref{thm:strict-convexity} proves global strict convexity in logit coordinates. (d) The exact cyclic family exhibits the fractional band \(2<\gamma<R_3(\alpha)\).}
\label{fig:overview-multipanel}
\end{figure}
